\documentclass[12pt,oneside]{book}

% =====================================================
% PAQUETES
% =====================================================
\usepackage[utf8]{inputenc}
\usepackage[T1]{fontenc}
\usepackage[spanish,es-nodecimaldot]{babel}

\usepackage[
    top=2.5cm,
    bottom=2.5cm,
    left=3cm,
    right=2.5cm,
    headheight=15pt
]{geometry}

\usepackage{amsmath}
\usepackage{amssymb}
\usepackage{mathtools}

\usepackage{graphicx}
\usepackage{float}
\usepackage{caption}
\usepackage{subcaption}

\usepackage{booktabs}
\usepackage{array}
\usepackage{longtable}
\usepackage{tabularx}

\usepackage{enumitem}

\usepackage{xcolor}
\usepackage[most]{tcolorbox}

\usepackage{fancyhdr}
\usepackage{ragged2e}

\graphicspath{
    {./img/}
    {./graficos/}
    {./figuras/}
}

% =====================================================
% METADATOS
% =====================================================
\newcommand{\universidad}{Universidad de Buenos Aires}
\newcommand{\facultad}{Facultad de Derecho}
\newcommand{\materia}{Derecho Civil --- Introducción al Derecho Privado}
\newcommand{\titulo}{Clasificación y Modalidades de los Actos Jurídicos}
\newcommand{\autor}{Isaac Edgar Camacho Ocampo}
\newcommand{\anio}{2025}

\usepackage[
    colorlinks=true,
    linkcolor=blue!50!black,
    citecolor=green!40!black,
    urlcolor=blue!60!black,
    pdftitle={\titulo},
    pdfauthor={\autor},
    pdfsubject={Apuntes universitarios},
    bookmarks=true
]{hyperref}

% =====================================================
% CONFIGURACIÓN GENERAL
% =====================================================
\setcounter{tocdepth}{2}

\setlength{\parindent}{0pt}
\setlength{\parskip}{0.5em}

\setlist[itemize]{
    topsep=0.3em,
    itemsep=0.2em
}
\setlist[enumerate]{
    topsep=0.3em,
    itemsep=0.2em
}

\clubpenalty=10000
\widowpenalty=10000

% =====================================================
% COMANDOS MATEMÁTICOS
% =====================================================
\newcommand{\abs}[1]{\left|#1\right|}
\newcommand{\norm}[1]{\left\lVert#1\right\rVert}
\newcommand{\vect}[1]{\mathbf{#1}}
\newcommand{\deriv}[2]{\frac{d#1}{d#2}}
\newcommand{\pderiv}[2]{\frac{\partial#1}{\partial#2}}

% =====================================================
% ENCABEZADOS Y PIES
% =====================================================
\pagestyle{fancy}
\fancyhf{}
\fancyhead[L]{\nouppercase{\leftmark}}
\fancyhead[R]{\materia}
\fancyfoot[C]{\thepage}
\renewcommand{\headrulewidth}{0.4pt}
\renewcommand{\footrulewidth}{0pt}

\fancypagestyle{plain}{
    \fancyhf{}
    \fancyfoot[C]{\thepage}
    \renewcommand{\headrulewidth}{0pt}
}

% =====================================================
% CAJAS ACADÉMICAS
% =====================================================
\newtcolorbox{definition}[1]{
    enhanced,
    breakable,
    title={Definición: #1},
    fonttitle=\bfseries,
    colback=blue!3,
    colframe=blue!50!black,
    boxrule=0.8pt,
    arc=2mm
}

\newtcolorbox{important}{
    enhanced,
    breakable,
    title={Importante},
    fonttitle=\bfseries,
    colback=yellow!8,
    colframe=orange!70!black,
    boxrule=0.8pt,
    arc=2mm
}

\newtcolorbox{example}[1]{
    enhanced,
    breakable,
    title={Ejemplo: #1},
    fonttitle=\bfseries,
    colback=green!4,
    colframe=green!40!black,
    boxrule=0.8pt,
    arc=2mm
}

\newtcolorbox{formula}[1]{
    enhanced,
    breakable,
    title={Fórmula / Regla: #1},
    fonttitle=\bfseries,
    colback=purple!4,
    colframe=purple!50!black,
    boxrule=0.8pt,
    arc=2mm
}

\newtcolorbox{exercise}[1]{
    enhanced,
    breakable,
    title={Ejercicio: #1},
    fonttitle=\bfseries,
    colback=gray!5,
    colframe=gray!60!black,
    boxrule=0.8pt,
    arc=2mm
}

\newtcolorbox{note}{
    enhanced,
    breakable,
    title={Nota},
    fonttitle=\bfseries,
    colback=white,
    colframe=gray!50,
    boxrule=0.6pt,
    arc=2mm
}

% =====================================================
% DOCUMENTO
% =====================================================
\begin{document}

% ---------------- PORTADA ----------------
\begin{titlepage}
\thispagestyle{empty}

\begin{center}

\vspace*{1cm}

{\large \textsc{\universidad}}

\vspace{0.2cm}

{\large \textsc{\facultad}}

\vspace{3cm}

{\Huge \textbf{\materia}}

\vspace{1cm}

{\LARGE \titulo}

\vspace{0.8cm}

\textit{Comprender la estructura de los actos jurídicos es el primer paso para razonar como jurista.}

\vspace{1.2cm}

\rule{0.70\textwidth}{0.5pt}

\vspace{0.5cm}

{\large Apuntes de estudio}

\vspace{0.5cm}

\rule{0.70\textwidth}{0.5pt}

\vfill

{\large \autor}

\vspace{0.3cm}

{\large \anio}

\end{center}

\vfill

\begin{flushleft}
\textit{Este es un modesto aporte a los estudiantes de ``Derecho Civil'' no pretende sustituir el material oficial de la materia.}
\end{flushleft}

\end{titlepage}

% ---------------- ÍNDICE ----------------
\tableofcontents

% =====================================================
\chapter{Clasificación de los Actos Jurídicos}
% =====================================================

En el régimen del Código de Vélez se enunciaba expresamente una serie de clasificaciones de los actos jurídicos --unilaterales y bilaterales, entre otras--. El Código Civil y Comercial de la Nación (CCyC) vigente, en cambio, optó metodológicamente por no incluir una clasificación sistematizada de los actos jurídicos: no existen artículos dedicados a definir de manera unificada las distintas clases de actos. Sin embargo, estas clasificaciones subsisten y se manifiestan a lo largo de todo el articulado, en la medida en que el Código atribuye consecuencias jurídicas distintas según el acto sea bilateral, entre vivos o de última voluntad, oneroso o gratuito, entre otras categorías.

\begin{note}
Estas clasificaciones no son puramente teóricas, sino que tienen consecuencias prácticas concretas y se encuentran presentes en distintos lugares del Código. A lo largo de las distintas asignaturas del plan de estudios de abogacía es habitual encontrar aplicaciones de estas categorías.
\end{note}

Su relevancia práctica es amplia: saber si un acto es formal o no formal determina si se requiere escritura pública; saber si es oneroso o gratuito modifica las reglas de buena fe aplicables (art. 392, CCyC); y saber si el acto está sujeto a condición, plazo o cargo indica en qué momento nace, se suspende o se extingue la obligación asumida. Estas categorías atraviesan los contratos, las sucesiones, el derecho de familia, las obligaciones y los derechos reales, y constituyen el marco conceptual necesario para interpretar el Código en su conjunto.

\begin{example}{Aplicación integradora de las clasificaciones}
En la redacción de una compraventa de un inmueble, dicho acto resulta simultáneamente bilateral, oneroso, formal (requiere escritura pública), patrimonial y de disposición. Si además las partes acuerdan que la entrega quede sujeta a que el banco apruebe el crédito del comprador, se configura una condición suspensiva; si se pacta que el contrato se resuelve en caso de que el banco no se expida dentro de los sesenta días, se configura un plazo resolutorio; y si el vendedor se compromete a pintar el departamento antes de la entrega, se configura un cargo. De este modo, las clasificaciones y las modalidades estudiadas en este material convergen en un único documento cotidiano de la práctica profesional.
\end{example}

\section{Actos Unilaterales y Bilaterales}

El criterio de esta clasificación es el número de partes que concurren a la celebración del acto.

\begin{definition}{Actos unilaterales y bilaterales}
El acto es \textbf{unilateral} cuando una sola parte concurre a su celebración. Es \textbf{bilateral} cuando dos o más partes concurren a su celebración.
\end{definition}

\begin{table}[H]
\centering
\begin{tabularx}{\textwidth}{l X l}
\toprule
\textbf{Tipo de acto} & \textbf{Definición} & \textbf{Ejemplo típico} \\
\midrule
Unilateral & Una sola parte concurre a la celebración del acto. & El testamento \\
Bilateral & Dos o más partes concurren a la celebración del acto. & Un contrato; el matrimonio \\
\bottomrule
\end{tabularx}
\caption{Actos unilaterales y bilaterales}
\end{table}

\section{Actos entre Vivos y de Última Voluntad}

\begin{definition}{Actos entre vivos y de última voluntad}
Los actos \textbf{entre vivos} son otorgados por personas vivas para producir efectos en vida de los otorgantes. Los actos \textbf{de última voluntad} son aquellos que producen sus efectos luego de la muerte de quien los otorga.
\end{definition}

\begin{table}[H]
\centering
\begin{tabularx}{\textwidth}{l X X}
\toprule
\textbf{Tipo de acto} & \textbf{Definición} & \textbf{Cuándo produce efectos} \\
\midrule
Entre vivos & Son otorgados por personas vivas para producir efectos en vida. & Desde el momento de su celebración \\
De última voluntad & Producen sus efectos luego de la muerte del otorgante. & Desde la muerte del otorgante \\
\bottomrule
\end{tabularx}
\caption{Actos entre vivos y de última voluntad}
\end{table}

\section{Actos a Título Oneroso y a Título Gratuito}

Esta clasificación resulta especialmente relevante a lo largo de todo el Código, dado que el CCyC tiende a proteger de manera diferenciada a quienes adquieren a título oneroso y de buena fe frente a quienes adquieren a título gratuito.

\begin{definition}{Actos onerosos y gratuitos}
El acto es \textbf{oneroso} cuando existen prestaciones recíprocas: cada parte entrega algo a cambio de lo que recibe. El acto es \textbf{gratuito} (o a título de liberalidad) cuando una de las partes recibe una prestación sin dar nada a cambio, lo cual enriquece a esa parte.
\end{definition}

\begin{table}[H]
\centering
\begin{tabularx}{\textwidth}{l X X}
\toprule
\textbf{Tipo de acto} & \textbf{Definición} & \textbf{Ejemplo} \\
\midrule
Oneroso & Hay prestaciones recíprocas: cada parte da algo a cambio. & Compraventa \\
Gratuito (liberalidad) & Una de las partes recibe sin dar nada a cambio; enriquece a esa parte. & Donación; renuncia a un derecho \\
\bottomrule
\end{tabularx}
\caption{Actos a título oneroso y a título gratuito}
\end{table}

El caso típico del título gratuito es la donación, pero también puede configurarse a título gratuito la renuncia a un derecho: quien tiene derecho a cobrar y renuncia a ese cobro está beneficiando a la otra parte sin recibir nada a cambio.

\begin{formula}{Art. 392, CCyC --- Efectos respecto de terceros en cosas registrables}
El Código diferencia la protección al subadquirente según haya adquirido a título oneroso o gratuito. El subadquirente de buena fe y a título oneroso goza de mayor protección.
\end{formula}

\section{Actos Principales y Accesorios}

\begin{definition}{Actos principales y accesorios}
Los actos \textbf{principales} son aquellos que no dependen, para su existencia, de otro acto. Los actos \textbf{accesorios} son aquellos que siguen la suerte del acto principal, del cual dependen para existir.
\end{definition}

\begin{table}[H]
\centering
\begin{tabularx}{\textwidth}{l X X}
\toprule
\textbf{Tipo de acto} & \textbf{Definición} & \textbf{Ejemplo} \\
\midrule
Principal & No depende de otro acto para existir. & Contrato de alquiler \\
Accesorio & Sigue la suerte del principal; depende de él para existir. & La garantía del contrato de alquiler \\
\bottomrule
\end{tabularx}
\caption{Actos principales y accesorios}
\end{table}

A modo de ejemplo, una garantía que acompaña a un contrato de alquiler constituye un acto accesorio respecto del principal y sigue su suerte, sin perjuicio de los matices particulares que se estudian en cada contrato en concreto.

\section{Actos Formales y No Formales}

Esta es una clasificación de particular importancia, que se retoma en profundidad al estudiar la forma del acto jurídico.

\begin{definition}{Actos formales y no formales}
Los actos \textbf{formales} son aquellos en los que la ley prescribe cierta solemnidad. Dentro de estos, la solemnidad puede hacer a la validez del acto (\textbf{formal solemne} o absolutamente formal) o hacer únicamente a su prueba (\textbf{formal no solemne} o relativamente formal). Los actos \textbf{no formales} son aquellos respecto de los cuales la ley no exige una forma determinada, quedando las partes en libertad de elegir la que prefieran.
\end{definition}

\begin{table}[H]
\centering
\begin{tabularx}{\textwidth}{l X X}
\toprule
\textbf{Tipo de acto} & \textbf{Definición} & \textbf{Consecuencia del incumplimiento de la forma} \\
\midrule
Formal solemne (absolutamente formal) & La ley prescribe una solemnidad cuyo incumplimiento invalida el acto. & Nulidad del acto \\
Formal no solemne (relativamente formal) & La ley prescribe una forma, pero su incumplimiento afecta solo la prueba, no la validez. & Puede probarse por otro medio \\
No formal & La ley no manda ninguna forma determinada; las partes pueden elegir la que prefieran. & No hay consecuencia: el acto es válido \\
\bottomrule
\end{tabularx}
\caption{Actos formales y no formales}
\end{table}

\begin{example}{Acto formal solemne}
El matrimonio es un acto formal solemne: se encuentra rodeado de una serie de solemnidades marcadas por el Código Civil y no puede celebrarse sino de acuerdo con los rituales que este establece. Cualquier otra forma determina que el matrimonio no sea válido.
\end{example}

\begin{example}{Acto no formal}
El alquiler de un inmueble es, en principio, un acto no formal: la ley no exige que se celebre por escritura pública. No obstante, las partes pueden optar por una forma más solemne que la exigida legalmente --por ejemplo, contratando a un escribano-- por razones de seguridad jurídica, sin que ello sea un requisito legal.
\end{example}

\section{Actos Patrimoniales y Extrapatrimoniales}

\begin{definition}{Actos patrimoniales y extrapatrimoniales}
Los actos \textbf{patrimoniales} tienen contenido económico y están vinculados con el patrimonio. Los actos \textbf{extrapatrimoniales} no tienen contenido económico directo y se vinculan con derechos personalísimos o con el derecho de familia.
\end{definition}

\begin{table}[H]
\centering
\begin{tabularx}{\textwidth}{l X X}
\toprule
\textbf{Tipo de acto} & \textbf{Definición} & \textbf{Ejemplo} \\
\midrule
Patrimonial & Tiene contenido económico; está vinculado con el patrimonio. & Un contrato de compraventa \\
Extrapatrimonial & No tiene contenido económico directo; vinculado con derechos personalísimos o de familia. & El reconocimiento de un hijo \\
\bottomrule
\end{tabularx}
\caption{Actos patrimoniales y extrapatrimoniales}
\end{table}

El reconocimiento de un hijo es un acto jurídico unilateral mediante el cual una persona expresa su voluntad de reconocerlo. Se trata de un acto típico del derecho de familia, extrapatrimonial en sí mismo, aun cuando pueda tener consecuencias patrimoniales indirectas, como el derecho a alimentos. Los actos vinculados con derechos personalísimos también son extrapatrimoniales.

\section{Actos de Disposición, de Administración y de Conservación}

\begin{definition}{Actos de disposición, administración y conservación}
Los actos de \textbf{disposición} producen una variación sustancial en la composición del patrimonio. Los actos de \textbf{administración} no producen un cambio sustancial en el patrimonio. Los actos de \textbf{conservación} están dedicados al mantenimiento de los bienes.
\end{definition}

\begin{table}[H]
\centering
\begin{tabularx}{\textwidth}{l X X}
\toprule
\textbf{Tipo de acto} & \textbf{Definición} & \textbf{Ejemplo} \\
\midrule
Disposición & Se produce una variación sustancial en la composición del patrimonio. & Venta de un inmueble: sale el bien, entra dinero \\
Administración & No produce un cambio sustancial en el patrimonio. & Cobrar un alquiler, recibir una renta \\
Conservación & Dedicados al mantenimiento de los bienes. & Reparaciones, actos de preservación \\
\bottomrule
\end{tabularx}
\caption{Actos de disposición, administración y conservación}
\end{table}

\begin{note}
Esta clasificación resulta especialmente relevante en el derecho de familia y de las personas, en materia de tutela, curatela y régimen de poderes. El representante legal de una persona incapaz, por ejemplo, necesita autorización judicial para los actos de disposición, pero no para los de administración ordinaria.
\end{note}

\section{Actos Puros y Simples, y Actos Modales}

\begin{definition}{Actos puros y simples, y actos modales}
Los actos \textbf{puros y simples} no están sujetos a ninguna modalidad y producen sus efectos de inmediato. Los actos \textbf{modales} son aquellos en los cuales las partes eligen sujetarlo a alguno de los tres modos posibles: condición, plazo o cargo.
\end{definition}

\begin{table}[H]
\centering
\begin{tabularx}{\textwidth}{l X}
\toprule
\textbf{Tipo de acto} & \textbf{Definición} \\
\midrule
Puro y simple & No está sujeto a ninguna modalidad. Produce efectos de inmediato. \\
Modal & Las partes eligen sujetarlo a alguno de los tres modos: condición, plazo o cargo. \\
\bottomrule
\end{tabularx}
\caption{Actos puros y simples, y actos modales}
\end{table}

\section{Actos Causales y Abstractos}

\begin{definition}{Actos causales y abstractos}
El acto \textbf{causal} permanece vinculado a su causa, la cual influye en su validez y en sus efectos. El acto \textbf{abstracto} se independiza de su causa.
\end{definition}

\begin{table}[H]
\centering
\begin{tabularx}{\textwidth}{l X}
\toprule
\textbf{Tipo de acto} & \textbf{Definición} \\
\midrule
Causal & El acto permanece vinculado a su causa; la causa influye en su validez y efectos. \\
Abstracto & El acto se independiza de su causa. \\
\bottomrule
\end{tabularx}
\caption{Actos causales y abstractos}
\end{table}

% TODO: contenido incompleto o ambiguo en las notas originales -- el apunte remite a un desarrollo previo de los actos abstractos ("se estudiaron previamente en la clase") que no forma parte de este documento.

\section{Actos Positivos y Negativos}

\begin{definition}{Actos positivos y negativos}
El acto \textbf{positivo} involucra una acción, es decir, un hacer. El acto \textbf{negativo} involucra una omisión, es decir, un no hacer.
\end{definition}

\begin{table}[H]
\centering
\begin{tabularx}{\textwidth}{l X X}
\toprule
\textbf{Tipo de acto} & \textbf{Definición} & \textbf{Ejemplo} \\
\midrule
Positivo & Involucra una acción (un hacer). & Entregar una cosa, pagar una suma \\
Negativo & Involucra una omisión (un no hacer). & Comprometerse a no competir en un rubro comercial \\
\bottomrule
\end{tabularx}
\caption{Actos positivos y negativos}
\end{table}

\begin{example}{Acto negativo: cláusula de no competencia}
Es habitual que quienes se dedican a una determinada actividad --por ejemplo, la producción de helados-- al vender la marca y la empresa se comprometan a no generar una nueva empresa ni emprendimientos vinculados con ese rubro durante los próximos cinco años. Este acto jurídico se encuadra dentro de los actos negativos, en tanto las partes se obligan a un ``no hacer'', es decir, a una omisión.
\end{example}

% =====================================================
\chapter{Modalidades del Acto Jurídico}
% =====================================================

Las modalidades del acto jurídico son la condición, el plazo y el cargo. En el régimen del Código de Vélez, las modalidades estaban reguladas dentro de la sección de obligaciones del Libro Segundo. En el Código Civil y Comercial vigente, en cambio, se las incorporó al Libro Primero, Título Cuarto, dentro del capítulo de hechos y actos jurídicos, en un capítulo específico dedicado a las modalidades del acto jurídico.

\section{La Condición}

\begin{formula}{Art. 343, CCyC --- Alcances}
Se denomina condición a la cláusula de los actos jurídicos por la cual las partes subordinan su plena eficacia o resolución a un hecho futuro e incierto.
\end{formula}

\begin{definition}{Condición}
La condición es una cláusula por la cual se subordina la adquisición o la resolución de un acto a un hecho futuro e incierto. Lo esencial de la noción es la incertidumbre: el hecho condicionante puede o no llegar a suceder.
\end{definition}

\begin{example}{Condición suspensiva (1)}
``Te voy a regalar un auto si Argentina sale campeón del mundo en el próximo mundial.'' Se trata de un hecho futuro e incierto: puede ocurrir o no. El efecto --la donación del auto-- queda suspendido hasta que ocurra, o no, el hecho condicionante.
\end{example}

\begin{example}{Condición suspensiva (2)}
``Te voy a regalar mi biblioteca jurídica si te recibís de abogado/a.'' El hecho es incierto: puede ocurrir o no. La obligación existe desde la constitución del acto, pero no es exigible --no tiene plena eficacia-- hasta el momento en que ocurre el hecho condicionante.
\end{example}

\begin{example}{Condición resolutoria}
``Te dono mi auto, pero siempre y cuando mi hermano no vuelva de España. Si mi hermano vuelve, tenés que devolverme el auto.'' Aquí el acto ya produce efectos desde el principio; el hecho condicionante produce la resolución, es decir, la devolución de los efectos hacia atrás.
\end{example}

\begin{table}[H]
\centering
\begin{tabularx}{\textwidth}{l X X}
\toprule
\textbf{Tipo de condición} & \textbf{Punto de partida} & \textbf{Efecto del hecho condicionante} \\
\midrule
Suspensiva & El acto existe pero NO produce efectos todavía. & Si ocurre el hecho, el acto produce sus efectos; si no ocurre, el acto no produce efectos. \\
Resolutoria & El acto YA produce efectos desde su celebración. & Si ocurre el hecho, los efectos cesan (se resuelven), con retroactividad. \\
\bottomrule
\end{tabularx}
\caption{Diferencia entre condición suspensiva y resolutoria}
\end{table}

\subsection{La Condición Impropia o Suposición}

El nuevo Código incorporó también la llamada ``condición impropia'' o ``suposición'': aquella en la que las partes condicionan el acto a un hecho presente o pasado que, al momento del acto, es ignorado por ellas.

\begin{formula}{Art. 343, CCyC (segundo párrafo)}
Se equipara a la condición la cláusula por la cual las partes sujetan la eficacia del acto a la existencia de cierto estado de cosas al tiempo de celebrarlo.
\end{formula}

\begin{example}{Condición impropia o suposición}
Una persona ignora la situación patrimonial de un familiar que vive lejos y, sin saberlo, celebra un acto sujeto a que ese familiar no haya caído en quiebra: por ejemplo, ``te presto la plata en la medida en que tal familiar mío no haya terminado en quiebra''. Aunque el hecho ya ocurrió o no ocurrió al momento de celebrar el acto, la incertidumbre subjetiva de las partes --que lo ignoran-- hace que este supuesto también valga como condición.
\end{example}

\subsection{Clasificación de las Condiciones}

\begin{table}[H]
\centering
\begin{tabularx}{\textwidth}{l l X}
\toprule
\textbf{Criterio} & \textbf{Tipo} & \textbf{Descripción} \\
\midrule
Según el hecho & Positiva & Implica una acción (que algo ocurra). \\
Según el hecho & Negativa & Implica una omisión (que algo no ocurra). \\
Según el efecto & Suspensiva & Subordina la eficacia del acto a un hecho futuro e incierto. \\
Según el efecto & Resolutoria & Subordina la extinción del acto a un hecho futuro e incierto. \\
Según la fuente & Causal & El hecho no depende de las partes; depende del azar o de un tercero. \\
Según la fuente & Potestativa & Depende de la voluntad del deudor o del acreedor. \\
Según la fuente & Mixta & Depende en parte del deudor y en parte de un tercero o del azar. \\
\bottomrule
\end{tabularx}
\caption{Clasificación de las condiciones}
\end{table}

La condición \textbf{causal} no depende de las partes, sino del azar o de un tercero --por ejemplo, que la selección nacional resulte campeona--. La condición \textbf{potestativa} depende del acreedor o del deudor --por ejemplo, ``te voy a regalar mi biblioteca si te recibís de abogado/a'' es potestativa del acreedor--. La condición \textbf{mixta} depende en parte del deudor y en parte de un hecho ajeno a él.

\begin{important}
La condición puramente potestativa del deudor está prohibida: ``te comprometo a entregarte algo si yo decido hacerlo'' no puede depender completamente de la voluntad del obligado.
\end{important}

\subsection{Condiciones Prohibidas}

\begin{formula}{Art. 344, CCyC --- Condiciones prohibidas}
Es nula la condición que: (1) sea un hecho imposible, (2) sea contraria a las buenas costumbres, (3) esté prohibida por el ordenamiento jurídico, (4) dependa exclusivamente de la voluntad del obligado (puramente potestativa). La condición de no hacer algo imposible no perjudica la validez de la obligación.
\end{formula}

\begin{formula}{Art. 345, CCyC --- Condiciones que gravan libertades fundamentales}
Las condiciones cuyo cumplimiento depende de que el beneficiario cambie de domicilio, religión u otras libertades fundamentales se tienen por no escritas.
\end{formula}

\begin{example}{Condición que grava una libertad fundamental}
``Te dono la casa si cambiás de religión'' constituye una condición prohibida, en tanto grava la libertad fundamental de profesar la religión que la persona elija.
\end{example}

\begin{table}[H]
\centering
\begin{tabularx}{\textwidth}{X X}
\toprule
\textbf{Tipo de condición prohibida} & \textbf{Consecuencia jurídica} \\
\midrule
Hecho imposible / contraria a buenas costumbres / puramente potestativa (art. 344, párr. 1) & Nulidad absoluta del acto \\
Que grava libertades fundamentales (art. 345, párr. 3) & Se tiene por no escrita (la obligación principal podría subsistir) \\
Imposible en testamentos (excepción) & La condición se tiene por no puesta; la disposición vale \\
\bottomrule
\end{tabularx}
\caption{Condiciones prohibidas y sus consecuencias}
\end{table}

Existe una diferencia entre el primer y el tercer párrafo del artículo 344: las condiciones ilícitas del primer párrafo producen la nulidad absoluta del acto, mientras que las condiciones comprendidas en el tercer párrafo se tienen por no escritas, discutiéndose en la doctrina si la obligación principal subsiste o no. La excepción se presenta en los testamentos: las condiciones prohibidas incorporadas en un testamento no afectan la validez de las disposiciones testamentarias de las que forman parte.

\section{El Plazo}

A diferencia de la condición, el plazo es la modalidad que posterga o limita temporalmente los efectos del acto. En el plazo hay un hecho futuro pero \textbf{cierto}: se sabe que va a ocurrir, aunque no siempre se sepa exactamente cuándo.

\begin{example}{Plazo suspensivo y plazo resolutorio}
Plazo suspensivo: ``Te voy a donar estas casas dentro de un año.'' No hay incertidumbre sobre si llegará el día, aunque pueda haberla sobre el momento exacto en el caso del plazo incierto. Plazo resolutorio: ``Te presto el uso de la casa por un año.'' Los efectos comienzan de inmediato, pero cesan cuando vence el plazo.
\end{example}

\begin{table}[H]
\centering
\begin{tabularx}{\textwidth}{l X X}
\toprule
\textbf{Tipo de plazo} & \textbf{Definición} & \textbf{Ejemplo} \\
\midrule
Suspensivo & Posterga los efectos del acto hasta el vencimiento del plazo. & ``Te dono la casa dentro de un año.'' \\
Resolutorio & Los efectos del acto cesan al vencimiento del plazo. & ``Te presto la casa por un año.'' \\
Cierto & Se conoce la fecha exacta del vencimiento. & ``En 30 días'', ``el 31 de diciembre'' \\
Incierto & No se sabe exactamente cuándo vencerá, pero es seguro que ocurrirá. & ``Cuando muera tal persona'' \\
Determinado & Las partes, la ley o el juez fijaron la fecha o el momento. & La ley o el contrato establece una fecha \\
Indeterminado & No se precisó fecha, pero existe un plazo posible. & Cláusula ``a mejor fortuna'' (art. 889, CCyC) \\
Esencial & El cumplimiento fuera de término no satisface el interés del acreedor. & Catering para una fiesta: si llega dos semanas tarde, ya no sirve \\
Accidental & El cumplimiento tardío puede igualmente satisfacer el interés del acreedor. & Entrega de materiales de construcción con demora \\
\bottomrule
\end{tabularx}
\caption{Clasificación del plazo}
\end{table}

\begin{example}{Plazo esencial}
Una persona contrata un servicio de catering para una fiesta y, por un error, el servicio se entrega dos semanas después del evento. El plazo era esencial --la fiesta misma-- por lo que la prestación tardía no satisface el interés del acreedor y el incumplimiento es total.
\end{example}

El plazo puede surgir de la convención de las partes, de la ley, puede ser judicial, o bien indeterminado, cuando se establece la fecha de inicio pero no la de finalización.

\begin{formula}{Art. 889, CCyC --- Cláusula ``a mejor fortuna''}
El Código admite la cláusula ``a mejor fortuna'': ``yo te presto plata y me devolvés cuando estés bien económicamente.'' Se trata de un plazo indeterminado. El juez, a solicitud del acreedor, puede fijar un plazo y hacer cumplir la obligación.
\end{formula}

\section{El Cargo}

\begin{definition}{Cargo}
El cargo es una obligación accesoria impuesta al adquirente de un derecho, que no impide los efectos del acto pero que, sin embargo, puede exigirse.
\end{definition}

\begin{example}{Cargo simple}
Una persona dona un campo a otra con el cargo de que mantenga en la casa a los caseros. Si el donatario no cumple, los caseros podrán exigir el cumplimiento del cargo o reclamar la indemnización de los daños y perjuicios ocasionados por el incumplimiento; sin embargo, en principio, el acto de donación no se cae.
\end{example}

\begin{example}{Cargo con finalidad benéfica}
Se dona a una fundación un terreno con el cargo de que el producto de la explotación económica de ese terreno se destine a obras para los barrios vulnerables. La fundación que recibe la donación debe administrar el bien conforme a ese cargo.
\end{example}

\begin{table}[H]
\centering
\begin{tabularx}{\textwidth}{X X X}
\toprule
\textbf{Característica} & \textbf{Condición} & \textbf{Cargo} \\
\midrule
Suspende efectos del acto & SÍ (si es suspensiva) & NO \\
Coercible (exigible) & Cuando se cumple el hecho: sí & Sí, siempre \\
Si falla en cumplirlo, revoca el acto & Sí (si es resolutoria) & Solo si fue impuesto como cargo condicional \\
Quién puede exigirlo & Las partes & Los beneficiarios y quien constituyó el derecho \\
\bottomrule
\end{tabularx}
\caption{Diferencias entre condición y cargo}
\end{table}

El cargo es coercitivo pero no suspensivo, a diferencia de la condición suspensiva, que sí suspende los efectos. El cargo siempre se impone al adquirente como una obligación que debe cumplir y, en tanto obligación, debe reunir los requisitos propios de las obligaciones: ser lícito, posible, determinado y de contenido patrimonial. Surge de la autonomía de la voluntad, es excepcional y accidental --puede o no estar presente--, constituye una obligación accesoria del derecho transmitido, debe ser expreso y es restrictivo del derecho adquirido, en la medida en que recibir algo con el cargo de cumplir una obligación restringe el derecho de quien lo recibe.

\subsection{Cargo Simple y Cargo Condicional}

\begin{table}[H]
\centering
\begin{tabularx}{\textwidth}{l X X}
\toprule
\textbf{Tipo de cargo} & \textbf{Definición} & \textbf{Consecuencia del incumplimiento} \\
\midrule
Cargo simple & Obligación accesoria que no condiciona la adquisición del derecho. & Se exige el cumplimiento o se reclaman daños. El acto no se revoca. \\
Cargo condicional & El cumplimiento del cargo fue puesto como condición del acto. & Si no se cumple, se puede revertir la adquisición del derecho. Opera como condición resolutoria. \\
\bottomrule
\end{tabularx}
\caption{Cargo simple y cargo condicional}
\end{table}

El cargo simple no tiene efecto resolutorio sobre el derecho principal. El cargo condicional, en cambio, opera como una condición resolutoria (o suspensiva), de modo que su incumplimiento afecta a la propia adquisición del derecho: si el cargo fue puesto como condicional, la adquisición puede revertirse en caso de incumplimiento.

\subsection{Cargos Prohibidos}

\begin{formula}{Art. 357, CCyC --- Cargos prohibidos}
Son de aplicación a los cargos las normas establecidas para las condiciones prohibidas (art. 344, CCyC). Sin embargo, cuando se impone un cargo prohibido, no se produce la nulidad del acto, sino que el cargo se tiene por no escrito.
\end{formula}

Los cargos prohibidos son, en su contenido, los mismos supuestos previstos para las condiciones prohibidas: el artículo 357 remite al artículo 344. La diferencia radica en la consecuencia: cuando se impone un cargo prohibido, este se tiene por no escrito, sin que ello acarree la nulidad del acto principal.

% =====================================================
\chapter{Cierre y Síntesis}
% =====================================================

Al finalizar el desarrollo de las clasificaciones y las modalidades del acto jurídico --condición, plazo y cargo-- corresponde una aclaración fundamental sobre la naturaleza de estas últimas.

\begin{important}
Las modalidades --condición, plazo y cargo-- \textbf{no} son una clasificación de los actos jurídicos. Son modos, cláusulas, elementos accidentales de un acto jurídico: pueden o no estar presentes. Si las partes las pactan, rigen las reglas de su regulación civil; si no las incorporan, el acto es puro y simple y produce efectos de inmediato.
\end{important}

Los actos jurídicos son manifestaciones de voluntad orientadas a producir efectos jurídicos. El Código Civil y Comercial no los clasifica en un artículo único, sino que sus categorías atraviesan todo el texto legal. Estas clasificaciones --unilateral/bilateral, entre vivos/última voluntad, oneroso/gratuito, principal/accesorio, formal/no formal, patrimonial/extrapatrimonial, disposición/administración/conservación, puro/modal, causal/abstracto, positivo/negativo-- no son puramente teóricas: determinan las reglas aplicables en cada situación jurídica concreta. Sobre esa base, el acto puede quedar sujeto a modalidades accidentales: la condición (hecho futuro e incierto que suspende o resuelve los efectos), el plazo (hecho futuro cierto que posterga o limita los efectos en el tiempo) y el cargo (obligación accesoria que no suspende el acto, pero que es coercitiva). Cada modalidad presenta variantes propias, reglas de validez específicas y artículos particulares dentro del CCyC (arts. 343 a 357).

% =====================================================
\chapter{Cuadro Cornell de Repaso}
% =====================================================

El siguiente cuadro reúne, según el método Cornell, las preguntas centrales que permiten repasar activamente el contenido desarrollado en los capítulos anteriores, junto con sus respuestas y una síntesis final integradora.

\begin{longtable}{
    >{\raggedright\arraybackslash}p{0.30\textwidth}
    >{\raggedright\arraybackslash}X
}
\toprule
\textbf{Preguntas / palabras clave} & \textbf{Notas / respuestas} \\
\midrule
\endfirsthead
\toprule
\textbf{Preguntas / palabras clave} & \textbf{Notas / respuestas} \\
\midrule
\endhead

\textbf{¿Qué son las clasificaciones de los actos jurídicos y para qué sirven?} &
Son categorías doctrinarias --y en gran parte también legales-- que agrupan los actos según criterios como el número de partes, el momento de producción de efectos, la existencia de prestaciones recíprocas, la dependencia respecto de otro acto, la forma, el contenido patrimonial o el nivel de variación del patrimonio. No son puramente teóricas: determinan las reglas aplicables al acto en cada situación concreta del Código y de la práctica profesional. \\

\textbf{Diferencia entre acto unilateral y bilateral. Ejemplos.} &
Unilateral: una sola parte (el testamento). Bilateral: dos o más partes (el contrato, el matrimonio). Esta diferencia resulta clave porque el régimen de perfeccionamiento, revocación y efectos es distinto en cada caso. \\

\textbf{¿Por qué importa si un acto es oneroso o gratuito?} &
El CCyC protege de manera distinta a los subadquirentes según hayan adquirido a título oneroso o gratuito. Conforme el art. 392, CCyC, quien adquiere de buena fe y a título oneroso goza de mayor protección frente a la evicción y otras acciones de terceros que quien adquirió gratuitamente. \\

\textbf{Acto formal solemne vs.\ formal no solemne vs.\ no formal.} &
Formal solemne: la forma es requerida para la validez; si no se cumple, el acto es nulo (testamento, matrimonio). Formal no solemne: la forma prescrita afecta solo la prueba; el acto igual es válido. No formal: la ley no exige ninguna forma; las partes eligen cómo exteriorizar su voluntad. \\

\textbf{¿Qué es la condición? Definición legal.} &
Cláusula por la cual se subordina la adquisición o la resolución de un acto a un hecho futuro e incierto (art. 343, CCyC). Lo esencial es la incertidumbre: el hecho puede o no ocurrir. Si el hecho es cierto, se trata de un plazo, no de una condición. \\

\textbf{Condición suspensiva vs.\ resolutoria.} &
Suspensiva: los efectos del acto no se producen hasta que ocurra el hecho condicionante; el acto existe, pero no es exigible. Resolutoria: el acto ya produce efectos; si ocurre el hecho, los efectos cesan y se resuelven retroactivamente. Ejemplo de suspensiva: ``te dono la biblioteca si te recibís''. Ejemplo de resolutoria: ``te dono el auto, pero si vuelve mi hermano, me lo devolvés''. \\

\textbf{¿Qué condiciones están prohibidas? Consecuencias.} &
Son nulas las condiciones que: (1) son imposibles, (2) son contrarias a las buenas costumbres, (3) son puramente potestativas del deudor (art. 344, CCyC). Se tienen por no escritas las que gravan libertades fundamentales, como la religión o el domicilio (art. 345). En los testamentos, las condiciones prohibidas no afectan la validez de la disposición testamentaria. \\

\textbf{Diferencia entre condición y plazo.} &
En ambos casos el hecho es futuro. La diferencia clave es la certeza: en la condición el hecho es incierto (puede no ocurrir); en el plazo el hecho es cierto (se sabe que ocurrirá, aunque no siempre se sepa exactamente cuándo, como en el plazo incierto: ``cuando muera tal persona''). \\

\textbf{Tipos de plazo según su certeza y determinación.} &
Cierto: se conoce la fecha exacta. Incierto: no se sabe cuándo, pero es seguro que ocurrirá. Determinado: fijado por las partes, la ley o el juez. Indeterminado: sin fecha precisada (por ejemplo, la cláusula ``a mejor fortuna'', art. 889, CCyC). Esencial: el cumplimiento tardío no satisface el interés del acreedor (catering para una fiesta). Accidental: el cumplimiento tardío igualmente puede satisfacer el interés del acreedor. \\

\textbf{¿Qué es el cargo? ¿Suspende los efectos del acto?} &
El cargo es una obligación accesoria impuesta al adquirente de un derecho que NO suspende los efectos del acto: el acto produce efectos de inmediato. Sin embargo, el cargo es coercitivo: puede exigirse su cumplimiento. Si fue impuesto como condición (cargo condicional), su incumplimiento puede revertir la adquisición del derecho. Los cargos prohibidos se tienen por no escritos (art. 357, que remite al art. 344, CCyC), sin que ello invalide el acto principal. \\

\textbf{¿Son las modalidades una clasificación de los actos jurídicos?} &
No. Las modalidades (condición, plazo, cargo) son elementos accidentales del acto jurídico: pueden o no estar presentes. Si las partes las incorporan, el acto queda sujeto a esas reglas; si no las incorporan, el acto es puro y simple y produce efectos de inmediato, sin condicionamiento alguno. \\

\midrule
\multicolumn{2}{p{\dimexpr\textwidth-2\tabcolsep}}{
\textbf{Resumen:} Los actos jurídicos admiten múltiples clasificaciones --según el número de partes, el momento de sus efectos, la existencia de contraprestación, su dependencia de otro acto, su forma, su contenido patrimonial y su incidencia en el patrimonio-- que determinan las reglas aplicables en cada caso concreto del Código Civil y Comercial. Además de estas clasificaciones, el acto jurídico puede quedar sujeto a modalidades accidentales: la condición (hecho futuro e incierto), el plazo (hecho futuro cierto) y el cargo (obligación accesoria y coercitiva que no suspende el acto). Cada modalidad tiene reglas propias de validez, variantes específicas y condiciones o cargos expresamente prohibidos por la ley (arts. 343 a 357, CCyC). Comprender estas categorías es la base del razonamiento jurídico: permiten leer el Código, redactar documentos, asesorar a clientes y resolver casos con precisión técnica en cualquier rama del derecho privado.
} \\
\bottomrule
\end{longtable}

\end{document}
